\section{Scale、Emergence 与 Alignment：能力从哪里来}

理解 attention 后，课程把视角从单层机制拉到大规模训练。这里有三条不能混为一谈的线：规模怎样改变可观测能力，alignment 怎样改变输出偏好，产品模型怎样在数据、架构和训练配方上形成不同取舍。

\subsection{Large Language Model 是被放大的训练系统}

Large language model（大语言模型，LLM）通常仍以 next-token prediction 为核心，但参数、数据与计算规模显著增加。它的“large”不仅指参数，还意味着数据 pipeline、并行训练、评估与部署基础设施共同扩大。

\lecturefigure{slide-028.jpg}{LLM 的规模、训练数据与 next-token objective}{官方 CS25 V4 slide deck，第 28 页}

\begin{importantbox}{Autoregressive objective}
对序列 $x_{1:T}$，训练损失常写成
\[
\mathcal{L}_{\mathrm{LM}}(\theta)=-\sum_{t=1}^{T}\log p_{\theta}(x_t\mid x_{<t}).
\]
$\theta$ 是模型参数，$x_{<t}$ 是历史 token。目标只要求提高下一个 token 的似然；遵循指令、引用证据、拒绝危险操作和长期规划都不是这个目标直接保证的能力。
\end{importantbox}

\begin{warningbox}{参数量不是能力的充分统计量}
同样参数规模下，token 数、数据质量、architecture、optimizer、context length、post-training 与评估污染都会改变结果。课堂后面讨论 Phi、BabyLM、DPO 和 MoE，正是在说明“更大”之外还有大量设计自由度。
\end{warningbox}

\subsection{Emergent ability 与度量陷阱}

Emergent ability（涌现能力）被定义为大模型出现而小模型没有的能力。课堂图展示若干任务在规模阈值附近突然上升，但讲者也引用研究质疑：离散或非线性 metric 可能把平滑改善显示成跳变。接下来的图要分两遍读：第一遍看任务曲线，第二遍检查纵轴定义与评分 threshold。

\lecturefigure{slide-029.jpg}{多项任务随模型规模出现表面突变}{官方 CS25 V4 slide deck，第 29 页}

\begin{knowledgebox}{读图：先分清横轴、纵轴与 threshold}
横轴通常是模型规模或训练计算，纵轴是 task metric。若 metric 只有“完全正确”才得分，底层概率的小幅连续改善也可能跨过阈值后突然记为 1。图能证明某个评估上的非线性，不能单独证明模型内部发生了离散相变。
\end{knowledgebox}

\lecturefigure{slide-030.jpg}{Few-shot prompting 随规模变化的任务表现}{官方 CS25 V4 slide deck，第 30 页}

\begin{importantbox}{Metric-induced emergence 的简单例子}
设模型给正确答案的概率随规模 $s$ 平滑增加为 $p(s)$。若评估要求连续 $k$ 个步骤全部正确，则整体准确率近似
\[
A(s)=p(s)^k.
\]
即使 $p(s)$ 平滑，$A(s)$ 也会在高区间急剧上升。更稳妥的分析应同时看 token likelihood、partial credit、calibration 与不同阈值下的曲线。
\end{importantbox}

\lecturefigure{slide-031.jpg}{对 emergent ability 的多种可能解释}{官方 CS25 V4 slide deck，第 31 页}

\teachervoice{讲者明确提出：emergence 可能部分来自研究者选择的非线性 metric，而不是模型能力真正突然出现。这不是否认 scale 的作用，而是要求把观测工具也纳入因果分析。}

\lecturefigure{slide-032.jpg}{Beyond Scaling：架构、数据与训练过程也能改变能力}{官方 CS25 V4 slide deck，第 32 页}

\begin{knowledgebox}{为什么 beyond scaling 值得单独讨论}
更好的 architecture 可以改善归纳偏置；更高质量数据提高每个 token 的学习价值；更稳定的 optimizer 与 curriculum 改善训练效率；retrieval、tool 和 memory 把部分能力移出参数。它们共同决定能力密度与可访问性。
\end{knowledgebox}

\lecturefigure{slide-033.jpg}{课堂向学生提出的 scale、民主化与 retrieval 问题}{官方 CS25 V4 slide deck，第 33 页}

\begin{warningbox}{课堂问题没有唯一答案}
“继续扩大模型是否足够”“参数记忆与 retrieval 哪个更好”都是设计轴，不是二选一。规模可带来通用表示，retrieval 提供可更新证据，小模型可降低成本，专用结构可提高效率。实际系统通常组合这些路线。
\end{warningbox}

\subsection{RLHF 与 DPO：从 likelihood 到 preference}

预训练让模型拟合语料分布，post-training 则让输出更接近人类偏好和任务规范。Reinforcement Learning from Human Feedback（基于人类反馈的强化学习，RLHF）常训练 reward model，再用强化学习优化 policy；Direct Preference Optimization（直接偏好优化，DPO）直接从偏好对构造分类式目标。

\lecturefigure{slide-035.jpg}{RLHF 的 preference collection、reward model 与 policy optimization}{官方 CS25 V4 slide deck，第 35 页}

\begin{importantbox}{Preference data 的基本形式}
对 prompt $x$，人类或规则比较两个回答 $y_w$ 与 $y_l$，其中 $y_w$ 是 preferred response。Reward model 常学习
\[
P(y_w\succ y_l\mid x)=\sigma\!\left(r_\phi(x,y_w)-r_\phi(x,y_l)\right).
\]
$r_\phi$ 是 reward score，$\sigma$ 是 logistic function。数据只表达相对偏好，不等于客观真理；标注者构成、rubric 与采样策略会进入模型行为。
\end{importantbox}

\lecturefigure{slide-036.jpg}{DPO 用 reference policy 直接优化 preference pair}{官方 CS25 V4 slide deck，第 36 页}

\begin{importantbox}{DPO objective 的直觉}
一种常见写法是
\[
\mathcal{L}_{\mathrm{DPO}}=-\mathbb{E}\log\sigma\!\left(\beta\left[
\log\frac{\pi_\theta(y_w\mid x)}{\pi_{\mathrm{ref}}(y_w\mid x)}-
\log\frac{\pi_\theta(y_l\mid x)}{\pi_{\mathrm{ref}}(y_l\mid x)}
\right]\right).
\]
$\pi_\theta$ 是待训练 policy，$\pi_{\mathrm{ref}}$ 是 reference model，$\beta$ 控制偏离 reference 的强度。DPO 简化了训练管线，但 preference bias、coverage 和 over-optimization 仍然存在。
\end{importantbox}

\teachervoice{课堂没有把 DPO 描述成“RLHF 已被淘汰”，而是强调 RLHF 依赖高质量反馈、reward 与 policy optimization，DPO 提供了一条更直接的 preference-learning 路线。}

\subsection{ChatGPT、GPT-4 与 Gemini：2024 年的课堂快照}

接下来的模型页属于 2024 年课程语境。它们用于说明 post-training、多模态、context、benchmark 与 MoE 等方向，不应被解读为 2026 年最新产品排名。阅读顺序应从训练阶段和 architecture 出发，再看当时 benchmark，而不是先依据品牌名称预设结论。

\lecturefigure{slide-037.jpg}{ChatGPT 与 GPT 系列的课堂概览}{官方 CS25 V4 slide deck，第 37 页}

\begin{knowledgebox}{读图：产品名称背后要拆成训练阶段}
一个 chat model 通常经历通用预训练、监督式 instruction tuning、preference optimization、安全数据与部署策略。界面相似不代表训练数据、模型结构或行为边界相同；公开信息不足时，不应从产品体验反推全部内部机制。
\end{knowledgebox}

\lecturefigure{slide-038.jpg}{GPT-4 的能力、评估与多模态课堂摘要}{官方 CS25 V4 slide deck，第 38 页}

\begin{warningbox}{Benchmark snapshot 不是永久排序}
分数依赖版本、prompt、评估集、污染程度与 judge。课堂给出的结果应保留日期，不能用旧 benchmark 声明今天某模型绝对领先。更重要的是理解评估维度：知识、reasoning、multimodality、safety 与 cost 并不由一个总分覆盖。
\end{warningbox}

\lecturefigure{slide-039.jpg}{Gemini 系列与 2024 年模型比较}{官方 CS25 V4 slide deck，第 39 页}

\lecturefigure{slide-040.jpg}{Gemini 课堂页用 MoE 解释 routed experts}{官方 CS25 V4 slide deck，第 40 页}

\begin{importantbox}{Mixture of Experts 的最小机制}
Mixture of Experts（专家混合，MoE）用 router 为 token 选择少量 expert：
\[
p_i(x)=\operatorname{softmax}(W_rx)_i,\qquad
h'=\sum_{i\in\operatorname{TopK}(p(x))}p_i(x)E_i(x).
\]
$W_r$ 是 router 参数，$p_i$ 是 expert 权重，$E_i$ 是第 $i$ 个 expert。Sparse activation 让总参数量增加时，每个 token 只计算少数 expert；代价是 load balancing、通信和 capacity 管理更复杂。
\end{importantbox}

\begin{warningbox}{“Expert”不等于可命名的人类专家}
某个 expert 可能呈现主题或功能偏好，但训练通常没有保证它只负责图像、数学或某种语言。把 MoE 直接画成稳定职业分工只是一种直觉，不能替代 routing 与 activation 的实测分析。
\end{warningbox}

\subsection{2024 的位置与下一步缺口}

课程用两页总结当时状态：LLM、alignment、多模态和 diffusion transformer 快速发展；下一步则包括更通用的 agent、效率、世界模型、持续学习与安全控制。本节用这组图承担章节转场：先区分已经形成产品的能力，再标出仍停留在研究目标的缺口，并为后面的跨域应用与系统限制建立问题清单。

\lecturefigure{slide-041.jpg}{Where we are：2024 年课堂对 LLM 生态的归纳}{官方 CS25 V4 slide deck，第 41 页}

\lecturefigure{slide-042.jpg}{The Future：更广泛应用、agent 与模型效率}{官方 CS25 V4 slide deck，第 42 页}

\begin{knowledgebox}{读图：Future 页混合了技术趋势与研究愿景}
有些方向已有系统原型，例如 multimodal generation、tool use 与 smaller models；另一些仍是目标，例如稳定自主 agent、低成本持续学习与更接近人类的数据效率。讲义会把“已经演示”与“可靠部署”分开。
\end{knowledgebox}

\lecturefigure{slide-043.jpg}{Missing ingredients：通向更通用系统的剩余问题}{官方 CS25 V4 slide deck，第 43 页}

\begin{importantbox}{能力扩展后的六个缺口}
\begin{enumerate}
  \item 更低的训练与推理成本；
  \item 对现实世界与因果关系更稳健的表示；
  \item 可更新的长期记忆与持续学习；
  \item 可检查、可编辑的内部机制；
  \item 长轨迹中的 error correction；
  \item 用户权限、security 与 value alignment。
\end{enumerate}
\end{importantbox}

\subsection{本章小结}

Scale 能扩大表示与任务覆盖，但 emergent curve 受 metric 影响；preference optimization 改变行为，却继承数据与目标偏差；MoE、retrieval、data quality 与 specialized training 提供 scale 之外的杠杆。模型页最有价值的部分不是旧排名，而是把能力拆成训练阶段和系统组件。

